\subsection{Maximal ordered fields}

DEFINITION 4. --- An ordered field K is maximal if every ordered algebraic extension of K is trivial.

Example. --- * We will see later (Gen.~Top., VIII, p.~1) that the field R of real numbers is a maximal ordered field. *

The existence of maximal ordered fields is a consequence of the following theorem:

THEOREM 2. --- Every ordered field K admits an ordered algebraic extension which is a maximal ordered field.

One can show that this ordered extension is unique up to K-isomorphism (VI, p.~40, Ex. 15).

Let \(\Omega\) be an algebraic closure of \(K\) , and let \(\mathfrak{N}\) be the set of pairs \((A, \omega)\) , where \(A\) is a sub- \(K\) -extension of \(\Omega\) , and \(\omega\) is an ordering on \(A\) making \(A\) an ordered extension of \(K\) . Order \(\mathfrak{N}\) by the relation « \(L\) is an ordered extension of \(M\) » between \(M\) and \(L\) . Equipped with this ordering, \(\mathfrak{N}\) is an inductive ordered set: indeed if \((L_{t})\) is a totally ordered family of elements of \(\mathfrak{N}\) , then the field \(L = \cup_{t} L_{t}\) , ordered by taking \(L_{+} = \cup_{t} (L_{t})_{+}\) is an upper bound for the

\(L_{i}\) . Then N has a maximal element, by Set Theory, III, p.~154, Th. 2, which completes the proof.

PROPOSITION 5. --- Let K be a maximal ordered field, and let f be a polynomial in K{[}X{]} which changes sign between two elements a and b of K (with a \textless{} b). Then f has a root x in K such that a \textless{} x \textless{} b.

At least one of the irreducible factors of f, say h, changes sign between a and b. Then the field \(\mathrm{K}[\mathrm{X}]/(h)\) admits the structure of an ordered extension of K (VI, p.~23, Prop. 3), and h has degree 1 (Def. 4). Since \(h(a)h(b)<0\) , the unique root x of h is such that a\textless x\textless b, since a polynomial function of degree 1 is monotonic.

PROPOSITION 6. --- Every positive element of a maximal ordered field K has a square root in K. Every polynomial of odd degree in K{[}X{]} has at least one root in K.

This follows immediately from Cor. 2 and 1 to Prop. 4 of VI, p.~24.

COROLLARY. --- On a maximal ordered field K there exists only one ordering compatible with the field structure.

Indeed the positive elements of \(\mathbf{K}\) are determined by its algebraic structure: they are the squares.

